% DAOconCoq.tex
% Demostracción asistida por ordenador con Coq
% José A. Alonso Jiménez <jalonso@us.es>
% Sevilla, 31 de Julio de 2018
% =============================================================================

\documentclass[a4paper,12pt,twoside]{book}

%%%%%%%%%%%%%%%%%%%%%%%%%%%%%%%%%%%%%%%%%%%%%%%%%%%%%%%%%%%%%%%%%%%%%%%%%%%%%%
%% § Paquetes adicionales                                                   %%
%%%%%%%%%%%%%%%%%%%%%%%%%%%%%%%%%%%%%%%%%%%%%%%%%%%%%%%%%%%%%%%%%%%%%%%%%%%%%%

% Configuración para XeLaTeX
\usepackage{fontspec}
\usepackage{xltxtra}
\defaultfontfeatures{Ligatures=TeX,Numbers=OldStyle}
\setromanfont{DejaVu Sans}
\setmonofont{DejaVu Sans Mono}[Scale={0.90}]

% Notas: La lista de fuentes disponibles se obtiene con fc-list

\usepackage[spanish]{babel}        % Castellanización.
\usepackage{fancyvrb}              % Verbatim extendido
\usepackage{a4wide}

\linespread{1.05}                  % Distancia entre líneas
\setlength{\parindent}{2em}        % Indentación de comienzo de párrafo
\raggedbottom                      % No ajusta los espacios verticales.

\usepackage[%
  colorlinks=true,
  urlcolor=blue,
  % pdftex,
  pdfauthor={José A. Alonso <jalonso@us.es>},%
  pdftitle={Demostración asistida por ordenador con Coq},%
  pdfstartview=FitH,%
  bookmarks=false]{hyperref}      

\setcounter{tocdepth}{1}
\setcounter{secnumdepth}{4}

\usepackage{tocstyle}
\usetocstyle{KOMAlike}

%%%%%%%%%%%%%%%%%%%%%%%%%%%%%%%%%%%%%%%%%%%%%%%%%%%%%%%%%%%%%%%%%%%%%%%%%%%%%%
%% § Cabeceras                                                              %%
%%%%%%%%%%%%%%%%%%%%%%%%%%%%%%%%%%%%%%%%%%%%%%%%%%%%%%%%%%%%%%%%%%%%%%%%%%%%%%

\usepackage{fancyhdr}

\addtolength{\headheight}{\baselineskip}

\pagestyle{fancy}

\cfoot{}

\fancyhead{}
\fancyhead[RE]{\mdseries\sffamily DAO con Coq}
\fancyhead[LO]{\mdseries\sffamily \nouppercase{\leftmark}}
\fancyhead[LE,RO]{\mdseries\sffamily \thepage}

%%%%%%%%%%%%%%%%%%%%%%%%%%%%%%%%%%%%%%%%%%%%%%%%%%%%%%%%%%%%%%%%%%%%%%%%%%%%%%
%% § Definiciones                                                           %%
%%%%%%%%%%%%%%%%%%%%%%%%%%%%%%%%%%%%%%%%%%%%%%%%%%%%%%%%%%%%%%%%%%%%%%%%%%%%%%

\input definiciones
\def\mtctitle{Contenido}

%%%%%%%%%%%%%%%%%%%%%%%%%%%%%%%%%%%%%%%%%%%%%%%%%%%%%%%%%%%%%%%%%%%%%%%%%%%%%%
%% § Título                                                                 %%
%%%%%%%%%%%%%%%%%%%%%%%%%%%%%%%%%%%%%%%%%%%%%%%%%%%%%%%%%%%%%%%%%%%%%%%%%%%%%%

\title{{\LARGE Demostración asistida por ordenador con Coq}}
\author{\href{http://www.cs.us.es/~jalonso}
        {\Large José A. Alonso Jiménez}}
\date{\vfill \hrule \vspace*{2mm}
  \begin{tabular}{l}
      \href{http://www.cs.us.es/glc}
           {Grupo de Lógica Computacional} \\
      \href{http://www.cs.us.es}
           {Dpto. de Ciencias de la Computación e Inteligencia Artificial} \\
      \href{http://www.us.es}
           {Universidad de Sevilla}  \\
      Sevilla, 31 de julio de 2018
  \end{tabular}\hfill\mbox{}}

%%%%%%%%%%%%%%%%%%%%%%%%%%%%%%%%%%%%%%%%%%%%%%%%%%%%%%%%%%%%%%%%%%%%%%%%%%%%%%
%% § Documento                                                              %%
%%%%%%%%%%%%%%%%%%%%%%%%%%%%%%%%%%%%%%%%%%%%%%%%%%%%%%%%%%%%%%%%%%%%%%%%%%%%%%

% \includeonly{recusion_sobre_numeros_naturales}

% \includexmp{licencia}

\begin{document}

\maketitle
\newpage

\noindent
Esta obra está bajo una licencia Reconocimiento--NoComercial--CompartirIgual
2.5 Spain de Creative Commons. 

\vspace*{4ex}

\begin{center}
\fbox{
\begin{minipage}{\textwidth}
\vspace*{1ex}
\noindent                     
{\bf Se permite:}
\begin{itemize}                         
\item copiar, distribuir y comunicar públicamente la obra
\item hacer obras derivadas
\end{itemize}
         
\noindent       
{\bf Bajo las condiciones siguientes:}

\begin{minipage}[t]{20mm}
  \vspace*{0mm}\hspace*{1em}
  \includegraphics[scale=0.35]{deed.png}
\end{minipage}%
\begin{minipage}[t]{14cm}
  \vspace*{0mm}
  {\bf Reconocimiento}. Debe reconocer los créditos de la  obra de la manera
  especificada por el autor.  
\end{minipage}
\\[1ex]

\begin{minipage}[t]{20mm}
  \vspace*{0mm}\hspace*{1em}
  \includegraphics[scale=0.35]{deed-eu.png}
\end{minipage}%
\begin{minipage}[t]{14cm}
  \vspace*{2mm}
  {\bf No comercial}. No puede utilizar esta obra para fines comerciales.
\end{minipage}
\\[1ex]

\begin{minipage}[t]{20mm}
  \vspace*{0mm}\hspace*{1em}
  \includegraphics[scale=0.35]{deed_002.png}
\end{minipage}%
\begin{minipage}[t]{14cm}
  \vspace*{-2mm}
  {\bf Compartir bajo la misma licencia}. Si altera o transforma esta obra, o
  genera una obra derivada, sólo puede distribuir la obra generada bajo una
  licencia idéntica a ésta. 
\end{minipage}

\begin{itemize}    
\item Al reutilizar o distribuir la obra, tiene que dejar bien claro los
  términos de la licencia de esta obra. 
\item Alguna de estas condiciones puede no aplicarse si se obtiene el permiso
 del titular de los derechos de autor.
\end{itemize}
\vspace*{1ex}
\end{minipage}}
\end{center}

\vspace*{4ex}

Esto es un resumen del texto legal (la licencia completa). Para ver una copia
de esta licencia, visite 
\href{http://creativecommons.org/licenses/by-nc-sa/2.5/es/}
     {\url{http://creativecommons.org/licenses/by-nc-sa/2.5/es/}}
o envie una carta a Creative Commons, 559 Nathan Abbott Way, Stanford,
California 94305, USA.

%%% Local Variables: 
%%% mode: latex
%%% TeX-master: t
%%% End: 

\newpage

\tableofcontents
\clearpage


\renewcommand{\chaptername}{Tema}

\chapter*{Introducción}

En este libro se incluye unos apuntes de demostración asistida por
ordenador con
\href{https://coq.inria.fr}
     {Coq}
para los cursos de

\begin{itemize}
\item
  \href{http://www.cs.us.es/~jalonso/cursos/m-ra}
       {Razonamiento automático}
  del
  \href{http://master.cs.us.es/Máster_Universitario_en_Lógica,_Computación_e_Inteligencia_Artificial}
       {Máster Universitario en Lógica, computación e inteligencia artificial}
  de la
  \href{http://www.us.es}
       {Universidad de Sevilla}.
\item 
  \href{http://www.cs.us.es/~jalonso/cursos/lmf}
       {Lógica matemática y fundamentos}
  del
  \href{http://www.us.es/estudios/grados/plan_171?p=7}
       {Grado en Matemáticas}
  de la
  \href{http://www.us.es}
       {Universidad de Sevilla}.
\end{itemize}

Esencialmente los apuntes son una adaptación del libro
\href{https://softwarefoundations.cis.upenn.edu/current/lf-current}
     {Software foundations (Vol. 1: Logical foundations)}
de Benjamin Peirce y otros.

Una primera versión de estos apuntes se han usado este año en el
\href{http://www.glc.us.es/~jalonso/SLC2018}
     {Seminario de Lógica Computacional}.

\chapter{Programación funcional y métodos elementales de demostración en Coq}
\teoria{T1_PF_en_Coq.v}

\chapter{ Demostraciones por inducción sobre los números naturales en Coq}
\teoria{T2_Induccion.v}

% \appendix
% \include{Examenes}

\end{document}

%%% Local Variables:
%%% mode: latex
%%% TeX-master: t
%%% TeX-engine: xetex
%%% End:
